\section{The power--log scale and its exact arithmetic}
\label{sec:scale}

\subsection{Blocks and their order}

Throughout, $w$ denotes a positive real variable tending to $0^+$, $L=\log w\to-\infty$ and $\M(w)=1+|\log w|$.

\begin{definition}[Blocks]
A \emph{power--log block} is an expression $w^\beta P(L)$ with $\beta\in\R$ and $P\in\R[L]$, $P\ne0$. Its \emph{exponent} is $\beta$ and its \emph{logarithmic degree} is $\deg P$.
\end{definition}

\begin{lemma}[Power beats logarithm]
\label{lem:dominance}
For every $\varepsilon>0$ and every real $b$, $w^\varepsilon\M(w)^b\to0$ as $w\to0^+$. Consequently, if $\beta<\gamma$ and $P,Q\ne0$, then $w^\gamma Q(L)=\oo\bigl(w^\beta P(L)\bigr)$; and within a fixed exponent, a higher logarithmic degree dominates: $w^\beta L^j=\oo(w^\beta L^k)$ for $j<k$.
\end{lemma}
\begin{proof}
Put $w=e^{-t}$, $t\to+\infty$: $w^\varepsilon\M(w)^b=e^{-\varepsilon t}(1+t)^b\to0$. For the second claim the ratio is $w^{\gamma-\beta}Q(L)/P(L)$, and $|Q(L)/P(L)|\le C\M(w)^{\deg Q}$ for $w$ small because $P(L)$ has no zeros for $L$ sufficiently negative and $|P(L)|\ge c|L|^{\deg P}\ge c$ there; the first claim finishes. The third claim is clear.
\end{proof}

The growth classes of nonzero blocks are totally ordered by increasing exponent and, at equal exponents, by decreasing logarithmic degree; blocks of the same exponent and degree have the same growth class. A finite sum of blocks with distinct exponents is asymptotic to its block of least exponent (whole polynomial included), and vanishes identically only if every block does. This is why truncations are taken by \emph{complete blocks}: cutting a block inside its polynomial would drop terms larger than the retained terms of the next exponent.

\begin{definition}[Remainder classes]
For $\beta\in\R$ and an integer $k\ge0$, let $\OO_{\beta,k}$ be the
class of functions $R$ satisfying
$|R(w)|\le Cw^\beta\M(w)^k$ for some $C>0$ and all sufficiently
small $w>0$. We use $R(w)=\OO_{\beta,k}$ as the usual shorthand for
membership in this class.
\end{definition}

By Lemma~\ref{lem:dominance}, $\OO_{\beta,k}\subseteq\OO_{\beta',k'}$ whenever $\beta>\beta'$, or $\beta=\beta'$ and $k\le k'$; the class with the smaller exponent, and at equal exponents the larger degree, is the coarser one. A block $w^\beta P(L)$ with $\deg P=k$ belongs to $\OO_{\beta,k}$ and to no $\OO_{\beta,k'}$ with $k'<k$. Sums and products of remainder classes obey
\begin{equation}
 \OO_{\beta,k}+\OO_{\beta',k'}=\OO_{\min(\beta,\beta'),\,k''},\qquad
 \OO_{\beta,k}\cdot\OO_{\beta',k'}=\OO_{\beta+\beta',\,k+k'},
 \label{eq:O-arithmetic}
\end{equation}
where $k''$ is the degree attached to the smaller exponent, and the larger of the two when the exponents coincide.

The order compares bounds on absolute magnitude, not signs of coefficients.
For example,
$\OO_{3,100}\subseteq\OO_{2,0}\subseteq\OO_{2,1}$:
the strict power advantage absorbs any fixed logarithmic degree, whereas at
the same exponent an extra logarithm weakens the bound. Combining two
independent errors uses the coarser class. Opposite signs in displayed terms
do not establish cancellation of their unknown errors.

\subsection{Exponent semigroups}

\begin{lemma}[Local finiteness]
\label{lem:finite}
Let $\delta_1,\dots,\delta_m>0$ and $S=\{\bk\cdot\bdelta:\bk\in\N^m\}$. For every real $H$ the set $\{\bk\in\N^m:\bk\cdot\bdelta\le H\}$ is finite, hence $S\cap(-\infty,H]$ is finite, and every $\sigma\in S$ has only finitely many representations $\sigma=\bk\cdot\bdelta$.
\end{lemma}
\begin{proof}
For $H<0$ the set is empty. For $H\ge0$, every admissible index satisfies
$0\le k_i\le\lfloor H/\delta_i\rfloor$ for each $i$, so there are at
most $\prod_i(1+\lfloor H/\delta_i\rfloor)$ indices. When $m=0$,
the sole index is the empty tuple and $S=\{0\}$.
\end{proof}

When the gaps are incommensurable the \emph{group} generated by them is dense in $\R$ (for instance $\Z+\sqrt2\Z$), but the semigroup of nonnegative combinations has only finitely many elements below any bound. This is the whole reason why expansions with irrational exponents can be computed term by term: the ordinary power--log expansions developed below have support in a translate $\beta_0+S$ of such a semigroup. Following \cite{proveit}, a subset of $\R$ that is finite below every bound is called \emph{left-finite}; the semigroups $S$ are left-finite, and so are their translates and finite unions.

The shift $\beta_0$ may be negative; positivity is required of the
\emph{gaps}, not of every exponent in the expansion. A negative generator
would give infinitely many exponents below a fixed bound, while a zero
generator would give infinitely many representations of the same weight.
Also, $\delta_*:=\min_i\delta_i$ bounds the increase along a single
generator step, not the spacing between distinct weights: for gaps $1$
and $\sqrt2$, the weights $1$ and $\sqrt2$ differ by less than
$\delta_*=1$. Local finiteness does not provide a fixed spacing for the
ordered support.

\subsection{Jets and their arithmetic}

\begin{definition}[Jets]
Fix gaps as above and $\beta_0\in\R$. A \emph{jet} (of cutoff $H$) is a finite sum $U=\sum_{\beta\in T}w^\beta C_\beta(L)$ with $T\subset\beta_0+S$ finite, $C_\beta\in\R[L]$, all $\beta<H$; equal exponents are merged and zero polynomials are discarded. Its \emph{valuation} $\val U$ is its least exponent ($+\infty$ for the empty jet).
Write $[F]_{<H}$ for the projection retaining complete blocks of exponent
less than $H$. For the algebra based at exponent zero, take finite sums
with exponents in $S\subseteq[0,\infty)$; its quotient $\PL_{<H}$,
$H>0$, uses addition and multiplication followed by this projection.
\end{definition}

In that nonnegative-exponent ring, finite sums supported at exponents
$\ge H$ form an ideal and $\PL_{<H}$ is the corresponding quotient.
A fixed translate $\beta_0+S$ need not itself be closed under products.
The quotient description also does not apply to the ring allowing arbitrary
real exponents: multiplication by a negative power can move a discarded term
below the cutoff. General shifted jets therefore require precision tracking;
a factor of valuation $\gamma<0$ lowers an unknown exponent cutoff by
$|\gamma|$ (Section~\ref{sec:evaluator}). Equal exponents must be identified
before their coefficient polynomials are added. In particular, exact
identities among algebraic exponents determine which terms belong to the
same block.

For instance, $[w^2]_{<2}=0$, but
$[w^{-1}w^2]_{<2}=w$. To multiply by $w^{-1}$ and determine the result
below exponent $2$, the original expression must be known below exponent
$3$. Likewise, an empty jet has infinite valuation as a finite
representative, but this says nothing about the valuation of an unknown
tail. Exact arithmetic on jets means exact determination of the retained
coefficients from the supplied data; it does not mean that the represented
function equals its finite jet.

This formal exponent filtration must be distinguished from an analytic
remainder class. For example, $w^H\log w$ is discarded at cutoff $H$,
but is not $\OO(w^H)$: its ratio to $w^H$ is unbounded. Thus a formal
cutoff alone cannot replace the pair $(H,k)$ in an analytic error bound.
The logarithmic degrees of retained blocks likewise do not determine
those of the tail. A complete-tail estimate requires the hypotheses of
Section~\ref{sec:evaluator} or another applicable asymptotic theorem.

\begin{lemma}[Reality of resolved coefficients]
\label{lem:real-resolved-coefficients}
Suppose $F$ is real-valued for all sufficiently small positive $w$, and
\[
 F(w)=\sum_{\beta\in B}w^\beta C_\beta(\log w)
       +\OO(w^P\M(w)^D),
\]
where $P\in\R$, $B$ is a finite set of distinct real exponents less than $P$,
$C_\beta\in\mathbb C[L]$, and $D\ge0$ is fixed and finite.
Then every $C_\beta$ has real coefficients.
\end{lemma}
\begin{proof}
Taking imaginary parts gives
$\sum_{\beta\in B}w^\beta\operatorname{Im}C_\beta(\log w)
=\OO(w^P\M^D)$, with imaginary parts taken coefficientwise.
If one of these real polynomials is nonzero, choose its least exponent
$\beta_0$. By Lemma~\ref{lem:dominance}, every higher block and the
remainder are of smaller order than
$w^{\beta_0}\operatorname{Im}C_{\beta_0}(\log w)$, a contradiction.
\end{proof}

The coefficients in this lemma are the complete polynomials after
equal-exponent collection; individual complex contributions may cancel.
The strict bound $\beta<P$ is needed: an imaginary term of order $w^P$
can be hidden in the error. Conversely, real coefficients in a finite
approximation do not prove that the original function is real-valued;
$1+i w^N=1+\OO(w^N)$ illustrates this for $N>0$.

\subsubsection*{Complex formal jets and successive expansions}
\label{sec:formal-successive-jets}

The same finite jet arithmetic can be defined with coefficient ring
$\mathbb C[L]$ in place of $\R[L]$, keeping the exponent support and
filtration unchanged. Here $L$ may be treated as an independent formal
symbol. Identities modulo higher powers assert algebraic consistency;
they do not select an analytic function, prove convergence, or supply an
analytic remainder bound. To interpret such coefficients on a complex
approach, a domain and the required branches of powers and logarithms
must be specified separately.

In a successive expansion in two variables, coefficients for the first
expansion may themselves be expanded in the second. The variable order
and the coefficient domains are part of this construction. Such an
iterated expansion does not by itself give a uniform bound as both
variables vary, nor justify substitution of a diagonal.

For example, for each fixed $a>0$,
\[
 F(w,a)=\frac{a}{a+w}=1+\OO_a(w)\qquad(w\to0^+),
 \qquad F(w,a)-1=-\frac{w}{a+w}.
\]
The subscript permits the bound constant to depend on $a$; indeed
$|F(w,a)-1|\le w/a$. On the diagonal $a=w$, however,
$F(w,w)=1/2$, so the conclusion $F(w,w)=1+\OO(w)$ is false.
Expanding instead in $a$ at fixed $w>0$ starts with $a/w$.
A simultaneous expansion or substitution along a varying parameter path
therefore needs a new argument or a uniform estimate on a domain
containing that path. This is the parameter-domain distinction used in
Section~\ref{sec:refinement-parameter-domain}.

\subsubsection*{Differentiation}

\begin{lemma}[Euler identity]
\label{lem:euler}
For $\beta\in\R$, $P\in\R[L]$ and $r\in\N$,
\begin{equation}
 \frac{\dd}{\dd w}\bigl[w^\beta P(L)\bigr]=w^{\beta-1}(\beta+\D)P(L),\qquad
 \frac{\dd^r}{\dd w^r}\bigl[w^\beta P(L)\bigr]=w^{\beta-r}\prod_{j=0}^{r-1}(\beta-j+\D)P(L).
 \label{eq:euler}
\end{equation}
The operators $\beta-j+\D$ commute. For nonzero $P$, $(\beta+\D)$
preserves $\deg P$ when $\beta\ne0$; when $\beta=0$ it lowers the
degree of a nonconstant polynomial by one and annihilates a constant.
\end{lemma}
\begin{proof}
$\frac{\dd}{\dd w}[w^\beta P(\log w)]=\beta w^{\beta-1}P+w^\beta P'(\log w)/w$. Iterate.
\end{proof}

The identity says that differentiation acts on a block as the scalar $\beta$ plus $\D$ on the polynomial; the exponent bookkeeping and the polynomial arithmetic separate completely. In the coefficient formulas below operators $\D+c$ appear with real $c$ (not necessarily positive), and their action on the coefficient vector $(z_0,\dots,z_d)$ of $P=\sum z_sL^s$ is bidiagonal: $(c+\D)P=\sum_s\bigl(cz_s+(s+1)z_{s+1}\bigr)L^s$.

\subsubsection*{Unit series}

Let $H>0$ and let $\Phi(t)=\sum_{k\ge0}\phi_kt^k$ be a formal power
series. At the zero jet, substitution gives $\Phi(0)=\phi_0$. For a
nonzero jet $U$ with $\nu:=\val U>0$, enlarge the positive generator
set to contain its finite support if necessary. Then $U^k$ has valuation
at least $k\nu$, so only finitely many powers survive in $\PL_{<H}$:
\begin{equation}
 \Phi(U)=\sum_{\substack{k\in\N\\k\nu<H}}\phi_kU^k\in\PL_{<H}.
 \label{eq:unit-series}
\end{equation}
In particular $(1+U)^a=\sum_k\binom akU^k$ for every real $a$ (with $\binom ak=a(a-1)\cdots(a-k+1)/k!$), $\log(1+U)=\sum_{k\ge1}(-1)^{k+1}U^k/k$, $e^U=\sum_kU^k/k!$, and $\phi(c+U)=\sum_k\phi^{(k)}(c)U^k/k!$ for any $\phi$ analytic at $c$. These are identities between finite sums of blocks modulo exponents $\ge H$. Substitution in a finite quotient does not require convergence of the original formal power series; analytic interpretation additionally uses the hypotheses of Sections~\ref{sec:forward} and \ref{sec:convergence}.

\subsubsection*{Composition with the logarithm}

The one operation that mixes the exponent bookkeeping with the polynomial coefficients is the substitution of a jet into the logarithm: if $w$ is replaced by $w(1+U)$ then $L$ becomes $L+\log(1+U)$, and a block $w^a P(L)$ becomes $w^a(1+U)^aP(L+\log(1+U))$. The following recurrence computes this without expanding $\log(1+U)$ separately.

\begin{lemma}[Composition recurrence]
\label{lem:composition}
Let $P\in\R[L]$ and $a\in\R$. For $|t|<1$,
\begin{equation}
 (1+t)^aP\bigl(L+\log(1+t)\bigr)=\sum_{k\ge0}Q_k(L)\,t^k,\qquad
 Q_0=P,\qquad Q_{k+1}=\frac{(a-k)Q_k+Q_k'}{k+1},
 \label{eq:composition}
\end{equation}
with $Q_k\in\R[L]$, $\deg Q_k\le\deg P$, and, when $a\in\N$ and $P$ is constant, $Q_k=0$ for $k>a$. More generally, if $Q_m$ is identically zero, then $Q_k=0$ for every $k\ge m$. Consequently, for a jet $U$ of positive valuation, $(1+U)^aP(L+\log(1+U))=\sum_kQ_k(L)U^k$ in $\PL_{<H}$, and an identically zero $Q_m$ ends this sum before any further powers of $U$ are needed.
\end{lemma}
\begin{proof}
Put $F(t)=(1+t)^aP(L+\log(1+t))$; it is analytic in $|t|<1$ with values in $\R[L]$, so $F=\sum Q_kt^k$ with $Q_k\in\R[L]$ and $Q_0=P$. Differentiating, $(1+t)\partial_tF=aF+\partial_LF$, because $\partial_t$ acting on the two factors gives $a(1+t)^{a-1}P+(1+t)^{a-1}P'$. Comparing coefficients of $t^k$ in $(1+t)\sum(k+1)Q_{k+1}t^k=\sum(aQ_k+Q_k')t^k$ gives $(k+1)Q_{k+1}+kQ_k=aQ_k+Q_k'$. The degree does not increase, and for $a\in\N$ and constant $P$ the factor $(a-k)$ annihilates $Q_{a+1}$ and hence all later $Q_k$ (for nonconstant $P$ the derivative keeps the sequence alive: $\log(1+t)$ is not a polynomial). In general, $Q_m=0$ implies $Q_m'=0$, so the homogeneous recurrence gives $Q_{m+1}=0$ and induction proves permanent annihilation. The jet statement follows because substitution of $U$ for $t$ is a ring homomorphism on finite sums.
\end{proof}

This stopping criterion uses the homogeneous recurrence, not merely an
observed zero coefficient of an arbitrary power series. For example,
$\sin t$ has zero even coefficients and nonzero odd coefficients. Likewise,
$Q_m(L_0)=0$ at a single value of $L$ does not imply $Q_m'=0$ and is
insufficient; the zero must be a polynomial identity under the fixed
parameter hypotheses.

\begin{remark}
The recurrence is the whole composition calculus needed here: the forward model $ax^p\bigl(1+\sum x^{\delta_i}B_i(\log x)\bigr)$ evaluated at $x=z(1+U)$ is
$az^p(1+U)^p\bigl(1+\sum_iz^{\delta_i}(1+U)^{\delta_i}B_i(L+\log(1+U))\bigr)$, each factor of which is an instance of \eqref{eq:composition} (with $P=1$ for the first). Its derivative is obtained from the same lemma with $a$ replaced by $a-1$ and $P$ by $aP+P'$, since $\frac{\dd}{\dd x}[x^aP(\log x)]=x^{a-1}(aP+P')(\log x)$.
\end{remark}

\subsection{Blocks with nonconstant leading coefficient}

Everything above allows the polynomial in the leading block to be nonconstant. In general, a real non-integer power or a logarithm of such a block produces $P(L)^a$ or $\log P(L)$ outside $\R[L]$. Special algebraic simplifications may stay in the class, but generic expressions require larger scales (rational functions of $L$, or iterated logarithms, Section~\ref{sec:boundaries}). Nonnegative integer powers of any jet stay in the class.
